Let $(n_k)_{k\ge 1}$ be a strictly increasing sequence of nonnegative integers and let $(a_k)_{k\ge 1}$ be nonzero complex numbers such that the power series
\[
f(z)=\sum_{k=1}^{\infty} a_k z^{n_k}
\]
converges for all $z\in\mathbb{C}$ (so $f$ is entire). Assume the gap condition
\[
\frac{n_k}{k}\to\infty\quad\text{as }k\to\infty.
\]
Question: Must it follow that $f$ assumes every complex value infinitely often (i.e., for every $w\in\mathbb{C}$, the equation $f(z)=w$ has infinitely many solutions in $\mathbb{C}$)?

Edge cases: $f$ is transcendental because infinitely many terms are present. By classical Picard theory, there is at most one value $w$ that can fail to be assumed infinitely often; the conjecture asks that there is no such exceptional value.
